\Article{8}
Werden Versuche dieser Art mit Magnetnadeln angestellt, die an einem senkrechten Faden aufgehängt sind, so darf bei genaueren Versuchen die Rückwirkung der Torsion nicht vernachlässigt werden. Wir unterscheiden an einem solchen Faden zwei horizontale Durchmesser. Der eine, $D$, liegt am unteren Ende, wo die Nadel befestigt ist, und verläuft parallel zur magnetischen Achse der Nadel. Der andere, $E$, liegt am oberen Ende, wo der Faden festgehalten wird, und verläuft bei unverdrehtem Faden parallel zu $D$. Wir nehmen an, dass $E$ mit dem magnetischen Meridian den Winkel $v$ bildet und die magnetische Achse oder $D$ mit ihm den Winkel $u$. Nach der Erfahrung ist die Torsionskraft wenigstens annähernd proportional zum Winkel $v-u$. Wir betrachten daher das Moment, mit dem diese Kraft den Winkel $u$ dem Winkel $v$ gleichzumachen sucht. Es ist $=(v-u)\theta$. Das Moment der erdmagnetischen Kraft, die den Winkel $u$ zu verkleinern sucht, ist nun $=TM\sin u$. Die Gleichgewichtsbedingung ist daher in der Gleichung $(v-u)\theta=TM\sin u$ enthalten. Diese lässt umso mehr reelle Lösungen zu, je kleiner $\theta$ im Verhältnis zu $TM$ ist. Da es hier nur um kleine Werte von $u$ geht, kann man sie ohne Bedenken durch $(v-u)\theta=TMu$ ersetzen, also durch $\frac vu=\frac{TM}{\theta}+1$. Bei unseren Apparaten ist das obere Ende des Fadens an einem beweglichen waagerechten Arm befestigt. Der Arm trägt einen Zeiger, der sich am Rand eines in Grade geteilten Kreises bewegt. Selbst wenn der Kollimationsfehler, also der Teilstrich, dem der Wert $v=0$ entspricht, noch nicht genau genug bekannt ist, zeigt der Zeiger den Unterschied zwischen zwei Werten von $v$ an. Ein anderer Teil des Apparats liefert ebenso mit größter Genauigkeit den Unterschied zwischen den zum Gleichgewicht gehörenden Werten von $u$. Offenbar erhält man $\frac{TM}{\theta}+1$, indem man den Unterschied zwischen zwei Werten von $v$ durch den Unterschied zwischen den entsprechenden Werten von $u$ teilt Verstreicht zwischen den hierzu angestellten Versuchen ein etwas längerer Zeitraum, so muss man, sofern höchste Genauigkeit gewünscht ist, die tägliche Änderung der magnetischen Deklination berücksichtigen. Das gelingt leicht mithilfe gleichzeitiger Beobachtungen an einem zweiten Apparat, bei dem das obere Fadenende unberührt bleibt. Es braucht kaum erwähnt zu werden, dass die beiden Apparate so weit voneinander entfernt sein müssen, dass sie sich nicht merklich stören können.

\expandafter\def\csname HistoricalNoteText9\endcsname{Genau genommen sind diese Einzelfäden keine wirklich einfachen Fäden. Es sind lediglich solche, die im Handel als ungesponnen verkauft werden.}
\SourceLeaf{98}
\expandafter\def\csname HistoricalNoteText8\endcsname{Auf beiden Skalen steigen die Teilstrichwerte von links nach rechts.}
Um zu zeigen, welche Genauigkeit derartige Beobachtungen zulassen, führen wir ein Beispiel aus dem Tagebuch an. Am $1832$ Sept $22$ beobachtete man bei unverändert belassenen Kollimationsfehlern die folgenden Deklinationen $u$ und Winkel $v$\HistoricalNote{8}:
\begin{center}
\begin{tabular}{c|c||c|c||c}
&&\multicolumn{2}{c||}{Erste Nadel}&Zweite Nadel\\
Vers.&Zeit&$u$&$v$&$u$\\\hline
I&$9^h33'$ vorm.&$+0^0\ 4'19''5$&$300^0$&$+0^0\ 2'12''1$\\
II&$9\ 57$&$-0\ 0\ 19{,}6$&$240$&$+0\ 1\ 37{,}7$\\
III&$10\ 16$&$-0\ 4\ 40{,}5$&$180$&$+0\ 1\ 18{,}8$
\end{tabular}
\end{center}
Die auf den Stand der ersten Beobachtung zurückgeführten Deklinationen der ersten Nadel sind somit
\begin{center}
\begin{tabular}{crr}
I.&$u=+0^0\ 4'19''5$&$v=300^0$\\
II&$+0\ 0\ 14{,}8$&$240$\\
III&$-0\ 3\ 47{,}2$&$180$
\end{tabular}
\end{center}
Aus der Verbindung der Beobachtungen ergeben sich für den Bruch $\frac{TM}{\theta}$ die folgenden Werte
\begin{center}
\begin{tabular}{cr}
I und II $.\ .\ .\ .\ .$&$881{,}7$\\
II und III $.\ .\ .\ .\ .$&$891{,}5$\\
I und III $.\ .\ .\ .\ .$&$886{,}6$
\end{tabular}
\end{center}

Die täglichen Änderungen der magnetischen Deklination werden durch die Torsion im Verhältnis der Einheit zu $\frac{n}{n+1}$ vermindert, wobei wir $\frac{TM}{\theta}=n$ setzen. Bei Fäden mit so geringer Torsion wie im vorangegangenen Beispiel kann diese Änderung als unmerklich gelten. Was die Dauer der unendlich kleinen Schwingungen betrifft, so lässt sich aus den Grundsätzen der Dynamik leicht folgern, dass sie durch die Torsion im Verhältnis der Einheit zu $\BareRoot{\frac{n}{n+1}}$ vermindert wird. Genau genommen gilt dies für den Fall $v=0$. Allgemein würden die Formeln gelten, wenn wir $\frac{TM\cos u^0}{\theta}=n$ setzten und mit $u^0$ den zum Gleichgewicht gehörenden Wert von $u$ bezeichneten. Der Unterschied wird aber ganz unmerklich sein.

\Article{9}
Der Koeffizient $\theta$ hängt hauptsächlich von der Länge, der Dicke und dem Material des Fadens ab. Bei Metallfäden hängt er außerdem etwas von der Temperatur ab, bei Seidenfäden von der Luftfeuchtigkeit. Dagegen scheint er bei jenen, vielleicht auch bei diesen, \textit{wenn sie einfach sind}, \SourceLeaf{99}überhaupt nicht von dem Gewicht abzuhängen, mit dem sie belastet werden. Anders verhält es sich bei zusammengesetzten Seidenfäden, wie man sie zum Tragen schwererer Nadeln verwenden muss. Bei diesen wächst $\theta$ mit dem angehängten Gewicht. Der Wert bleibt jedoch weit unter dem Wert von $\theta$ für einen Metallfaden gleicher Länge, der dasselbe Gewicht tragen kann. So wurde etwa mit einem Verfahren, das dem im vorigen Artikel beschriebenen ganz ähnlich ist, jedoch mit einem anderen Faden und einer anderen Nadel, der Wert $n=597{,}4$ gefunden, als der Faden die Nadel allein mit ihrem üblichen Zubehör trug und das Gesamtgewicht $496{,}2$ Gramm betrug. Dagegen ergab sich der Wert $=424{,}8$, als das Gewicht auf $710{,}8$ Gramm erhöht wurde. Somit war im ersten Fall $\theta=0{,}0016740\,TM$, im zweiten $\theta=0{,}0023542\,TM$. Der Faden ist $800$ Millimeter lang und aus $32$ einfachen Fäden\HistoricalNote{9} zusammengesetzt. Jeder trägt sicher etwa $30$ Gramm, und sie sind so angeordnet, dass sie alle unter gleicher Spannung stehen. Übrigens ist es wahrscheinlich, dass der Wert von $\theta$ aus einem konstanten und einem zum Gewicht proportionalen Teil besteht und dass der konstante Teil gleich der Summe der Werte von $\theta$ für die einzelnen einfachen Fäden ist. Unter dieser Hypothese, die bisher durch die Versuche noch nicht hinreichend bestätigt ist, ergibt sich für den konstanten Teil im angeführten Beispiel der Wert $=0{,}0001012\,TM$ und damit für den Wert von $\theta$ bei einem einfachen Faden $=0{,}00000316\,TM$. Mithilfe des gleich zu bestimmenden Werts von $TM$ lässt sich aus dieser Hypothese folgern, dass die Rückwirkung eines einfachen Fadens, der um einen dem Radius gleichen Bogen $(57^0\ 18')$ verdreht wird, der Gewichtskraft eines Milligramms entspricht, das auf einen Hebel von etwa $\frac1{17}$ Millimeter Länge drückt.

\Article{10}
Ist der schwingende Körper eine einfache Nadel von regelmäßiger Gestalt und homogener Masse, so lässt sich das Trägheitsmoment $K$ nach bekannten Verfahren berechnen. Ist der Körper zum Beispiel ein rechtwinkliges Parallelepiped mit den Seiten $a,b,c$, der Dichte $=d$ und somit der Masse $q=abcd$, so ist das Trägheitsmoment bezüglich einer durch die Mitte verlaufenden und zur Seite $c$ parallelen Achse $=\frac1{12}(aa+bb)q$. Bei Magnetnadeln dieser Gestalt ist die zur magnetischen Achse parallele Seite $a$ in der Regel viel länger als die Breite $b$, sodass es für weniger genaue Versuche sogar genügt, $K=\frac1{12}aaq$ zu setzen. Bei genaueren Versuchen würden wir jedoch auch bei Verwendung einer einfachen Nadel die unbegründete Annahme einer vollkommen homogenen Masse und einer vollkommen regelmäßigen Gestalt kaum \SourceLeaf{100}zulassen. Bei unseren Versuchen schwingt die Nadel mit einer aufwendigeren Ausstattung, sodass sich die Aufgabe durch eine solche Rechnung überhaupt nicht lösen lässt. Daher mussten wir uns ein anderes Verfahren überlegen, um das Moment $K$ mit größter Genauigkeit zu bestimmen.

Mit der Nadel war ein hölzerner Querstab verbunden, an dem zwei gleiche Gewichte hingen, die über sehr scharfe Spitzen auf die Punkte $A,B$ des Stabs drückten. Diese Punkte lagen auf einer waagerechten Geraden in derselben Vertikalebene wie die Aufhängungsachse und auf beiden Seiten in gleichem Abstand von ihr. Bezeichnet man die Masse jedes der beiden Gewichte mit $p$ und den Abstand $AB$ mit $2r$, so vergrößert sich das Moment $K$ durch Hinzufügen dieser Vorrichtung um $C+2prr$. Dabei ist $C$ die Summe aus dem Trägheitsmoment des Stabs bezüglich der Aufhängungslinie und den Momenten der Gewichte bezüglich der vertikalen Achsen durch ihre Spitzen und Schwerpunkte. Beobachtet man also die Schwingungen der unbelasteten Nadel sowie die der in zwei verschiedenen Abständen belasteten Nadel, etwa für $r=r'$ und $r=r''$, und erhält man für die Schwingungsdauern nach Zurückführung auf unendlich kleine Schwingungen und Ausschaltung der Wirkung der Torsion die Werte $=t,t',t''$ in dieser Reihenfolge, so lassen sich aus der Verbindung der Gleichungen
\[
\begin{aligned}
TMtt&=\pi\pi K\\
TMt't'&=\pi\pi(K+C+2pr'r')\\
TMt''t''&=\pi\pi(K+C+2pr''r'')
\end{aligned}
\]
die drei Unbekannten $TM,K$ und $C$ bestimmen. Noch größere Genauigkeit erreichen wir, wenn wir die Schwingungen für mehrere Werte von $r$, etwa $r=r',r'',r'''$ usw., mit den zugehörigen Schwingungsdauern $t',t'',t'''$ usw. beobachten und die beiden Unbekannten $x,y$ nach der Methode der kleinsten Quadrate so bestimmen, dass die folgenden Gleichungen möglichst genau erfüllt sind
\[
\begin{aligned}
t'&=\BareRoot{\frac{r'r'+y}{x}}\\
t''&=\BareRoot{\frac{r''r''+y}{x}}\\
t'''&=\BareRoot{\frac{r'''r'''+y}{x}}\ \text{usw.}
\end{aligned}
\]
Damit erhalten wir
\[
\begin{aligned}
TM&=2\pi\pi px\\
K+C&=2py
\end{aligned}
\]

\SourceLeaf{101}
Zu diesem Verfahren ist noch Folgendes anzumerken.

I. Verwendet man eine Nadel, die nicht allzu glatt poliert ist, so genügt es, den Holzstab einfach daraufzulegen. Ist die Oberfläche der Nadel aber so glatt, dass die Reibung ein Gleiten des Stabs darauf nicht verhindern kann, muss der Stab fester mit dem übrigen Apparat verbunden werden, damit sich der ganze Apparat wie ein einziger starrer Körper bewegt. In beiden Fällen ist darauf zu achten, dass die Punkte $A,B$ hinreichend genau auf einer waagerechten Geraden liegen.

II. Da eine Reihe solcher Versuche mehrere Stunden beansprucht, darf die Veränderlichkeit der Intensität des Erdmagnetismus in dieser Zeit nicht vernachlässigt werden, wenn höchste Genauigkeit gewünscht ist. Bevor man mit der Elimination beginnt, sind daher die beobachteten Schwingungsdauern auf einen konstanten Wert von $T$ zurückzuführen, etwa auf den Mittelwert für den ersten Versuch. Hierzu sind wie in Artikel 8 gleichzeitige Beobachtungen an einer anderen Nadel erforderlich. Ergeben diese als Dauer einer Schwingung für die mittleren Zeiten der einzelnen Versuche die Werte $=u,u',u'',u'''$ usw. in dieser Reihenfolge, so sind für die Rechnung anstelle der beobachteten Werte $t',t'',t'''$ usw. jeweils folgende Werte zu verwenden
\[
\frac{ut'}{u'},\qquad\frac{ut''}{u''},\qquad\frac{ut'''}{u'''}\quad\text{usw.}
\]

III. Ein entsprechender Hinweis gilt für die Veränderlichkeit von $M$, die sich aus einer etwaigen Temperaturänderung während der Versuche ergibt. Die eben beschriebene Zurückführung enthält diese Korrektur offensichtlich bereits, sofern beide Nadeln derselben Temperaturänderung unterworfen waren und von ihr in gleicher Weise beeinflusst werden.

IV. Wenn nur der Wert von $TM$ zu bestimmen ist, erübrigt sich der erste Versuch offenbar. Dennoch ist es nützlich, an die Versuche mit der belasteten Nadel unmittelbar einen Versuch mit der unbelasteten Nadel anzuschließen. Dadurch erhält man zugleich den Wert von $K$, der späteren Versuchen mit derselben Nadel zugrunde gelegt werden kann. Dieser Wert bleibt offenbar unverändert, auch wenn sich $T$ und $M$ im Laufe der Zeit ändern können.
